% Circuito equivalente c.c. da Questão 2 (base comum com PNP; capacitores abertos).
% Compilar com pdflatex; o PDF resultante substitui figuras/fig_q2_cc.pdf.
\documentclass[11pt,border=4pt]{standalone}
% Estilo comum dos esquemáticos do relatório (circuitikz).
% Uso: \documentclass[11pt,border=4pt]{standalone} e, em seguida, % Estilo comum dos esquemáticos do relatório (circuitikz).
% Uso: \documentclass[11pt,border=4pt]{standalone} e, em seguida, % Estilo comum dos esquemáticos do relatório (circuitikz).
% Uso: \documentclass[11pt,border=4pt]{standalone} e, em seguida, \input{estilo_circuitos}.
\usepackage[T1]{fontenc}
\usepackage[utf8]{inputenc}
\usepackage[brazilian]{babel}
\usepackage{lmodern}
\usepackage{amsmath}
\usepackage[american]{circuitikz}
\usetikzlibrary{calc,arrows.meta}

% Cores das figuras originais
\definecolor{anot}{RGB}{0,135,200}      % anotações: valores calculados, setas, resistências vistas
\definecolor{sombra}{RGB}{220,242,253}  % região do transistor nos modelos T e pi
\definecolor{term}{RGB}{62,124,148}     % nomes de terminais destacados (B, C, E)

% Fios com 0,4 pt e corpo dos componentes com 0,8 pt, como nas figuras originais
\ctikzset{bipoles/thickness=2}
.
\usepackage[T1]{fontenc}
\usepackage[utf8]{inputenc}
\usepackage[brazilian]{babel}
\usepackage{lmodern}
\usepackage{amsmath}
\usepackage[american]{circuitikz}
\usetikzlibrary{calc,arrows.meta}

% Cores das figuras originais
\definecolor{anot}{RGB}{0,135,200}      % anotações: valores calculados, setas, resistências vistas
\definecolor{sombra}{RGB}{220,242,253}  % região do transistor nos modelos T e pi
\definecolor{term}{RGB}{62,124,148}     % nomes de terminais destacados (B, C, E)

% Fios com 0,4 pt e corpo dos componentes com 0,8 pt, como nas figuras originais
\ctikzset{bipoles/thickness=2}
.
\usepackage[T1]{fontenc}
\usepackage[utf8]{inputenc}
\usepackage[brazilian]{babel}
\usepackage{lmodern}
\usepackage{amsmath}
\usepackage[american]{circuitikz}
\usetikzlibrary{calc,arrows.meta}

% Cores das figuras originais
\definecolor{anot}{RGB}{0,135,200}      % anotações: valores calculados, setas, resistências vistas
\definecolor{sombra}{RGB}{220,242,253}  % região do transistor nos modelos T e pi
\definecolor{term}{RGB}{62,124,148}     % nomes de terminais destacados (B, C, E)

% Fios com 0,4 pt e corpo dos componentes com 0,8 pt, como nas figuras originais
\ctikzset{bipoles/thickness=2}

\begin{document}
\begin{circuitikz}
  % Coordenadas em cm. Linha dos nós V_E e V_C em y = 0; transistor no meio.
  \def\xE{0}      % ramo de R_E e nó V_E
  \def\xQ{2.4}    % transistor (base para baixo)
  \def\xC{4.8}    % ramo de R_C e nó V_C

  % Transistor PNP: emissor à esquerda, coletor à direita, base para baixo, à terra
  \draw (\xQ,0) node[pnp, rotate=90] (Q) {};
  \draw (Q.B) node[ground] {};

  % Ramo do emissor: +20 V, R_E, corrente I_E e nó V_E
  \draw (\xE,2.72) to[R] (\xE,0) node[circ] {} -- (Q.E);
  \draw (\xE,2.72) node[vcc] {$+20\,\mathrm{V}$};
  \draw (\xE,0.40) node[currarrow, rotate=-90] {};

  % Ramo do coletor: nó V_C, corrente I_C, R_C e -25 V
  \draw (Q.C) -- (\xC,0) node[circ] {} to[R] (\xC,-2.72) node[vee] {$-25\,\mathrm{V}$};
  \draw (\xC,-0.40) node[currarrow, rotate=-90] {};

  % Rótulos
  \node[anchor=east, align=left, inner sep=0pt] at (\xE-0.40,1.36) {$R_E$\\$10\,\mathrm{k\Omega}$};
  \node[anchor=west, inner sep=0pt] at (\xE+0.26,0.40) {$I_E$};
  \node[anchor=north east, inner sep=0pt] at (\xE-0.12,-0.10) {$V_E$};
  \node[inner sep=0pt] at (\xQ,0.75) {$Q$};
  \node[anchor=south west, inner sep=0pt] at (\xC+0.16,0.20) {$V_C$};
  \node[anchor=east, inner sep=0pt] at (\xC-0.22,-0.40) {$I_C$};
  \node[anchor=west, align=left, inner sep=0pt] at (\xC+0.40,-1.36) {$R_C$\\$6{,}5\,\mathrm{k\Omega}$};
\end{circuitikz}
\end{document}
